\section{Software Evaluation}
\label{sec:sweval}

Before committing silicon area to the scheduling core of
Section~\ref{sec:arch}, the eviction and migration policies must be
validated on realistic agent traces. This section evaluates the
\unison policy along five dimensions that, taken together, stress
different failure modes of session-level caching.
Section~\ref{sec:sw-main} compares aggregate performance against
baseline and published policies on six traces.
Section~\ref{sec:sw-comp} ablates scoring and tiering to confirm their
complementarity. Section~\ref{sec:sw-signal} quantifies each signal
term via leave-one-out analysis on the full evaluation grid.
Section~\ref{sec:sw-robust} verifies that the gain persists under
storage-hierarchy perturbation, and Section~\ref{sec:sw-lut} examines
hazard lookup table (LUT) fairness and online convergence. Finally,
Section~\ref{sec:sw-vllm} replaces the trace-driven simulator with a
live vLLM server to validate that the policy transfers to a production
block-cache pool.

\subsection{Experimental Setup and Metrics}
\label{sec:sw-setup}

The evaluation replays six agent traces constructed by crossing the
SWE-bench~\cite{jimenez_swebench_2024} and
GAIA~\cite{mialon_gaia_2023} task suites with three model families,
namely Qwen3-Coder-30B~\cite{qwen3_2025},
Devstral-24B~\cite{devstral_2025}, and
Gemma4-E4B~\cite{gemma4_2026}, yielding 1\,415 sessions and 33\,596
turns in total. No public multi-agent serving trace with session-level
annotations exists as of this writing, so this single-agent replay
corpus is the most detailed open benchmark available. The main comparison includes recency (LRU), timeout,
the AGSERVE expected-completion estimator, CacheScout, \unison, and the
offline B\'el\'ady oracle. CacheScout is evaluated on a matched
two-tier harness at the Edge-Tight envelope against its own LRU
reference. Scoring-only and tiering-only variants are deferred to
Fig.~\ref{fig:ablation}. Six SRAM-and-HBM capacity envelopes span edge
through CXL configurations. Hit rate and prefill reduction are
arithmetic means over these envelopes, while AMAT/LRU and TTFT/LRU are
geometric-mean ratios relative to recency. TTFT is reported at the p50
percentile under $4\times$ load.

A request is an SRAM hit, an HBM hit, or a miss. With $R$ requests the
hit rate and the average memory access time are
\begin{equation}
  \mathrm{HR}=\frac{H_{\mathrm{S}}+H_{\mathrm{H}}}{R},
  \qquad
  \mathrm{AMAT}=\frac{1}{R}\sum_{i=1}^{R}
    \begin{cases}
      \ell_i\,t_{\mathrm{s}} & \text{SRAM hit},\\
      \ell_i\,t_{\mathrm{h}} & \text{HBM hit},\\
      \ell_i\,t_{\mathrm{p}} & \text{miss},
    \end{cases}
  \label{eq:hr-amat}
\end{equation}
where $t_{\mathrm{s}}$, $t_{\mathrm{h}}$, and $t_{\mathrm{p}}$ are the
per-token SRAM, HBM, and prefill costs. Prefill reduction is the
fraction of those tokens that a hit avoids,
\begin{equation}
  \mathrm{PreRED}=1-\frac{\sum_i \ell_i^{\mathrm{comp}}}{\sum_i \ell_i},
  \label{eq:prered}
\end{equation}
and the B\'el\'ady ratio is
\begin{equation}
  \mathrm{BR}
  =\frac{\mathrm{HR}}{\mathrm{HR}^{\star}}
  \label{eq:br}
\end{equation}
against the offline oracle. Prefetch accuracy is defined only when \tide
is present. Serving TTFT follows the same prefill accounting under a
trace-driven load model.

\subsection{Main Results}
\label{sec:sw-main}

Tab.~\ref{tab:main} compares \unison with baseline recency and timeout
rules, the AGSERVE expected-completion estimator, CacheScout, and the
offline B\'el\'ady oracle across all six traces. \unison is the best
non-oracle entry on hit rate and AMAT on every dataset, achieving a
B\'el\'ady ratio of $\mathrm{BR}{=}0.93$. It surpasses LRU on AMAT in
all 36 matched comparisons with a mean reduction of 34.8\% and
outperforms CacheScout on hit rate across all six datasets. The
remaining 7\% gap to B\'el\'ady reflects the inherent cost of
causality, since an online policy cannot observe the next arrival and
the residual is the portion that~\eqref{eq:esp} does not close.

\paragraph{TTFT and the hit-rate versus latency trade-off}
Although \unison dominates on hit rate and AMAT, the AGSERVE
expected-completion estimator achieves a marginally lower aggregate
TTFT/LRU ratio. The serving model charges leftover prefill tokens to
the request at the head of the GPU queue. The estimator evicts the
session it predicts will return latest and can therefore discard a
long-prefix session that the \unison completion penalty retains. The
retained prefixes improve hit rate and AMAT but leave a longer residual
prefill on the critical path of a competing arrival, allowing the
estimator to gain on first-token latency while losing on every other
metric. The
0.01-point TTFT difference therefore reflects a fundamental
hit-rate versus latency trade-off rather than a ranking deficiency.

More broadly, when long prefills saturate GPU compute as on
SWE/Qwen3 and SWE/Devstral, the bottleneck shifts to queuing delay and
policy ordering ceases to influence TTFT. The cache-bound regime where
eviction policy most directly improves latency is isolated in the vLLM
validation of Section~\ref{sec:sw-vllm}. Scoring and tiering ablations
are presented in Fig.~\ref{fig:ablation}.

% Auto-generated by paper/tables/gen_tab4_main.py
\begin{table*}[t]
\centering
\renewcommand{\arraystretch}{1.15}
\setlength{\tabcolsep}{2pt}
\scriptsize
\begin{threeparttable}
\caption{Overall algorithmic performance of representative session-level policies across six agent-trace families.\tnote{a}}
\label{tab:main}
\begin{tabular*}{\textwidth}{@{\extracolsep{\fill}}lrrrrrrrrrrrrrrr@{}}
\toprule
\textbf{Policy} & \multicolumn{2}{c}{\textbf{SWE/Qw3}} & \multicolumn{2}{c}{\textbf{SWE/Dev}} & \multicolumn{2}{c}{\textbf{SWE/Ge4}} & \multicolumn{2}{c}{\textbf{GAIA/Qw3}} & \multicolumn{2}{c}{\textbf{GAIA/Dev}} & \multicolumn{2}{c}{\textbf{GAIA/Ge4}} & \multicolumn{3}{c}{\textbf{Aggregate}\tnote{b}} \\
\cmidrule(lr){2-3} \cmidrule(lr){4-5} \cmidrule(lr){6-7} \cmidrule(lr){8-9} \cmidrule(lr){10-11} \cmidrule(lr){12-13} \cmidrule(lr){14-16}
 & \textbf{HR} & \textbf{AMAT} & \textbf{HR} & \textbf{AMAT} & \textbf{HR} & \textbf{AMAT} & \textbf{HR} & \textbf{AMAT} & \textbf{HR} & \textbf{AMAT} & \textbf{HR} & \textbf{AMAT} & \textbf{BR} & \textbf{PreRED} & \textbf{TTFT}\tnote{c} \\
 & (\%) & /LRU & (\%) & /LRU & (\%) & /LRU & (\%) & /LRU & (\%) & /LRU & (\%) & /LRU & & (\%) & /LRU \\
\midrule
Recency (LRU) & 20.7 & 1.00 & 15.2 & 1.00 & 48.3 & 1.00 & 72.8 & 1.00 & 83.3 & 1.00 & 63.0 & 1.00 & 0.63 & 55.7 & 1.00 \\
Timeout (TTL) & 21.6 & 0.95 & 16.2 & 0.97 & 48.7 & 0.99 & 73.4 & 1.01 & 83.6 & 0.96 & 65.0 & 0.98 & 0.64 & 56.5 & 0.95 \\
ETA (AGSERVE)~\cite{agserve_2025} & 32.5 & 0.78 & 26.8 & 0.82 & 56.8 & 0.87 & 77.9 & 0.83 & 83.4 & 0.99 & 72.7 & 0.75 & 0.80 & 63.4 & \textbf{0.63} \\
CacheScout~\cite{cachescout_2026} & 2.6 & 1.00 & 2.8 & 1.00 & 4.4 & 1.00 & 31.9 & 0.95 & 67.2 & 0.98 & 24.0 & 0.98 & -- & 26.0 & 0.97 \\
\rowcolor{gray!10}
\textbf{\textsc{Unison}} & \textbf{43.8} & \textbf{0.65} & \textbf{30.5} & \textbf{0.78} & \textbf{65.3} & \textbf{0.55} & \textbf{82.3} & \textbf{0.49} & \textbf{83.6} & \textbf{0.72} & \textbf{76.7} & \textbf{0.52} & \textbf{0.93} & \textbf{66.6} & 0.64 \\
\textcolor{gray}{B\'el\'ady oracle} & 43.9 & 0.68 & 39.9 & 0.68 & 65.3 & 0.76 & 81.6 & 0.69 & 83.6 & 0.95 & 77.2 & 0.67 & 1.00 & 67.8 & 0.43 \\
\bottomrule
\end{tabular*}
\begin{tablenotes}\scriptsize
\item[a] Qw3 stands for Qwen3-Coder-30B, Dev stands for Devstral-24B and Ge4 stands for Gemma4-E4B.
\item[b] Hit rate(HR) and prefill reduction(PreRED) are means over six capacity envelopes.
\item[c] Average memory access time(AMAT/LRU) and time to first token(TTFT/LRU) metrics are geomean ratios to recency, and TTFT is p50 at $4\times$ load.
\end{tablenotes}
\end{threeparttable}
\end{table*}


\subsection{Scoring and Tiering Ablation}
\label{sec:sw-comp}

Scoring and tiering could in principle be substitutes, so
Fig.~\ref{fig:ablation} crosses both evictors with and without \tide to
test this hypothesis. The joint \unison configuration outperforms either
component in isolation because scoring determines which sessions to
retain while tiering determines where they reside. A strong evictor
without migration strands hot sessions in HBM, and a migrator without
\esp promotes the wrong sessions, yielding a prefetch accuracy of only
1.1\%.

\begin{figure*}[t]
\centering
\begin{subfigure}{0.24\textwidth}
\includegraphics[width=\linewidth]{fig04a_hr.pdf}
\caption{Hit rate}
\label{fig:ablation-hr}
\end{subfigure}\hfill
\begin{subfigure}{0.24\textwidth}
\includegraphics[width=\linewidth]{fig04b_amat.pdf}
\caption{AMAT / LRU}
\label{fig:ablation-amat}
\end{subfigure}\hfill
\begin{subfigure}{0.24\textwidth}
\includegraphics[width=\linewidth]{fig04c_prered.pdf}
\caption{Prefill reduction}
\label{fig:ablation-prered}
\end{subfigure}\hfill
\begin{subfigure}{0.24\textwidth}
\includegraphics[width=\linewidth]{fig04d_pfacc.pdf}
\caption{Prefetch accuracy}
\label{fig:ablation-pfacc}
\end{subfigure}
\caption{Component ablation of scoring and tiering.}
\label{fig:ablation}
\end{figure*}

\subsection{Signal Contribution Analysis}
\label{sec:sw-signal}

The hardware scorer of Section~\ref{sec:quant} retains only two of the
three \spear signals, so a leave-one-out analysis on the full 36-point
grid is needed to validate that reduction.
Fig.~\ref{fig:signal} confirms that hazard and gap recurrence are both
indispensable, while elapsed idle contributes the least with a mean
impact of 1.55\,pp on hit rate. Because elapsed idle is also the only
term that requires a per-slot timer in hardware, its omission from the
register-transfer level (RTL) implementation is justified on both
accuracy and area grounds.

\begin{figure}[t]
\centering
\begin{subfigure}{0.48\linewidth}
\includegraphics[width=\linewidth]{fig05a_dhr.pdf}
\caption{$\Delta$HR vs full-3sig}
\label{fig:signal-hr}
\end{subfigure}\hfill
\begin{subfigure}{0.48\linewidth}
\includegraphics[width=\linewidth]{fig05b_damat.pdf}
\caption{$\Delta$AMAT vs full-3sig}
\label{fig:signal-amat}
\end{subfigure}
\caption{Mechanism-signal ablation of \spear via leave-one-out on the 36-point evaluation grid.}
\label{fig:signal}
\end{figure}

\subsection{Robustness Across Capacity Operating Points}
\label{sec:sw-robust}

The preceding results use a single default storage envelope.
Tab.~\ref{tab:robust} verifies that the advantage is not an artifact of
that operating point by sweeping three hierarchy axes, namely HBM capacity
with SRAM fixed, SRAM capacity with HBM fixed, and the SRAM-to-HBM
partition with the total held at 60K tokens. The six-trace mean remains
positive on hit rate, AMAT, and prefill reduction at every tabulated
point. The gain grows monotonically with the SRAM share because a
larger fast tier amplifies the benefit of accurate residency ranking.
It diminishes only when HBM is large enough that recency alone
approaches the hit-rate ceiling, a regime in which two GAIA traces tie
on hit rate while the six-trace mean remains positive. Ancillary
parameters, including read latency, DMA bandwidth, and prefill rate, do
not reorder the policies and shift the mean AMAT gain by less than 0.6
percentage points around the 19.6\% baseline, so they are omitted from
the table.

% Auto-generated by paper/tables/gen_tab6_robust.py
\begin{table}[t]
\centering
\renewcommand{\arraystretch}{1.12}
\setlength{\tabcolsep}{3pt}
\scriptsize
\begin{threeparttable}
\caption{Gain robustness of \unison against LRU under storage-hierarchy perturbation.}
\label{tab:robust}
\begin{tabular*}{\columnwidth}{@{\extracolsep{\fill}}lrrr@{}}
\toprule
 & \multicolumn{3}{c}{\textbf{Gain against LRU}\tnote{a}} \\
\cmidrule(lr){2-4}
\textbf{Point} & \textbf{$\Delta$HR}\tnote{b} & \textbf{$\Delta$AMAT} &
\textbf{$\Delta$Pre} \\
 & (\%) & (\%) & (\%) \\
\midrule
\rowcolor{gray!10}
\multicolumn{4}{c}{\textbf{HBM capacity, SRAM fixed at 10K tokens}} \\
30K & $+14.5$ & $+15.9$ & $+11.1$ \\
50K$^{\star}$ & $+15.5$ & $+19.6$ & $+12.4$ \\
100K & $+15.0$ & $+24.4$ & $+11.8$ \\
200K & $+10.4$ & $+16.2$ & $+5.3$ \\
\rowcolor{gray!10}
\multicolumn{4}{c}{\textbf{SRAM capacity, HBM fixed at 50K tokens}} \\
5K & $+13.0$ & $+15.9$ & $+8.7$ \\
10K$^{\star}$ & $+15.6$ & $+19.6$ & $+12.4$ \\
20K & $+18.4$ & $+26.5$ & $+15.2$ \\
40K & $+25.2$ & $+38.0$ & $+24.1$ \\
\rowcolor{gray!10}
\multicolumn{4}{c}{\textbf{SRAM:HBM partition, total fixed at 60K tokens}} \\
5K:55K & $+13.5$ & $+15.7$ & $+9.2$ \\
10K:50K$^{\star}$ & $+15.5$ & $+19.6$ & $+12.4$ \\
20K:40K & $+18.4$ & $+26.7$ & $+15.3$ \\
30K:30K & $+22.8$ & $+30.9$ & $+19.4$ \\
40K:20K & $+28.8$ & $+34.7$ & $+25.3$ \\
\bottomrule
\end{tabular*}
\begin{tablenotes}\scriptsize
\item[a] Each cell is the six-trace mean of the gain against LRU, with the starred row representing the Edge-Tight point.
\item[b] Hit-rate and prefill-reduction gains are percentage points, and the AMAT gain is the percent reduction in AMAT.
\end{tablenotes}
\end{threeparttable}
\end{table}


\subsection{Hazard LUT Fairness and Online Convergence}
\label{sec:sw-lut}

All experiments so far rely on a hazard LUT fitted on the full trace,
which could bias the evaluation by overfitting its own sessions or by
embedding deployment-time knowledge unavailable during online serving.

Tab.~\ref{tab:lutcv} answers the first concern with a five-fold
cross-validation on SWE/Qwen3. Because a scatter of the same 20
fold-by-capacity cells would cluster on $y{=}x$, the table reports the
residual. The held-out LUT deviates from the all-session LUT by at
most 0.85\,pp and by 0.17\,pp on average, confirming that the hazard
distribution is a stable property of the turn-survival law rather than
of a particular session split.

Fig.~\ref{fig:fair-online} addresses the second concern by measuring
how quickly an initially empty LUT converges to the offline
performance. Each trace incrementally rebuilds the table as sessions
complete, accumulating the counts in~\eqref{eq:hazard} with $N$
ranging from 128 to 496. To compare traces with different absolute
hit-rate scales on a common frame, we define the attainment metric
\begin{equation}
  A(n)=\frac{H_{\mathrm{on}}(n)-H_{\mathrm{LRU}}}
            {H_{\mathrm{off}}-H_{\mathrm{LRU}}}
  \label{eq:attainment}
\end{equation}
as the fraction of the offline-minus-LRU hit-rate surplus realized
after $n$ observed sessions, where $A{=}0$ corresponds to the LRU
floor and $A{=}1$ to the offline table. All six convergence curves
reach $A{=}1$. Residual fluctuations on GAIA/Gemma4 reflect count
noise from its 128-session corpus, and the two Devstral traces start
near the ceiling because LRU already approaches the offline hit rate
on those workloads. The globally fitted LUT is therefore an evaluation
convenience rather than a hidden prior, as a production deployment
accumulates the same table incrementally at runtime.

\begin{table}[t]
\centering
\renewcommand{\arraystretch}{1.12}
\setlength{\tabcolsep}{4pt}
\scriptsize
\begin{threeparttable}
\caption{Cross-validation evidence for hazard-mechanism generalizability on SWE/Qwen3.}
\label{tab:lutcv}
\begin{tabular*}{\columnwidth}{@{\extracolsep{\fill}}lrrrr@{}}
\toprule
\textbf{Capacity} & \textbf{All-sess.\tnote{a}} & \textbf{Held-out\tnote{a}} &
\textbf{Mean $\Delta$ \tnote{b}} & \textbf{Max $\lvert\Delta\rvert$} \\
\midrule
100K & 4.90 & 4.88 & $-$0.03 & 0.11 \\
200K & 18.94 & 18.82 & $-$0.11 & 0.82 \\
300K & 38.75 & 38.53 & $-$0.21 & 0.78 \\
500K & 63.53 & 63.20 & $-$0.33 & 0.85 \\
\bottomrule
\end{tabular*}
\begin{tablenotes}\scriptsize
\item[a] The reported metric is five-fold mean hit rate in percent.
\item[b] $\Delta$ is held-out minus all-session in percentage points.
\end{tablenotes}
\end{threeparttable}
\end{table}


\begin{figure}[t]
\centering
\includegraphics[width=\columnwidth]{fig06_lut_online.pdf}
\caption{Online convergence of the hazard mechanism.}
\label{fig:fair-online}
\end{figure}

\subsection{Production-Stack Validation with vLLM}
\label{sec:sw-vllm}

The preceding experiments use a trace-driven simulator that abstracts
away serving-stack scheduling and memory-management details. To
validate that the \spear ranking transfers to a real inference server,
we integrate the score into the vLLM~\cite{kwon_pagedattention_2023}
v1 prefix cache, replacing the default LRU victim order with the
\spear score derived exclusively from runtime-observable metadata.
Paired replays of the GAIA/Qwen3 trajectory issue token-identical loads
at 32-way concurrency and zero temperature. Tool-call frequency is
varied between tight and loose schedules to span the cache-contention
range that representative agent workloads exhibit. All pairings
complete without request errors, and Tab.~\ref{tab:vllm} reports the
measured gain against LRU under three deployment cases that each
isolate one operational factor.

% Auto-generated by paper/tables/gen_tab7_vllm.py
\begin{table}[t]
\centering
\renewcommand{\arraystretch}{1.12}
\setlength{\tabcolsep}{2pt}
\scriptsize
\begin{threeparttable}
\caption{Measured \spear gain under vLLM production-stack deployment.\tnote{a}}
\label{tab:vllm}
\begin{tabular*}{\columnwidth}{@{\extracolsep{\fill}}llllrrr@{}}
\toprule
\textbf{Case} & \textbf{Model} & \textbf{Pool} &
\textbf{Gap}\tnote{b} &
\textbf{$\Delta$HR} & \textbf{$\Delta$TTFT}\tnote{c} &
\textbf{$\Delta$e2e} \\
 & & & & (pp) & (\%) & (\%) \\
\midrule
Cache-bound & 1.5B & 60K & Tight & \textbf{+3.4} & $-35$ & $-72$ \\
Capacity-stressed & 1.5B & 30K & Loose & \textbf{+2.5} & $-6$ & $-4$ \\
Native-model & 30B & 60K & Tight & \textbf{+1.6} & $-1$ & $-3$ \\
\bottomrule
\end{tabular*}
\begin{tablenotes}\scriptsize
\item[a] Workload is GAIA/Qwen3 at 32-way concurrency.
1.5B is Qwen2.5-1.5B and 30B is the trajectory-native
Qwen3-Coder-30B-AWQ. Pool sizes are 60K and 30K tokens.
\item[b] Tight/Loose denotes the diverse tool-gap schedule.
\item[c] Deltas are versus paired LRU in vLLM and
TTFT is the mean and e2e is p99, in percent.
\end{tablenotes}
\end{threeparttable}
\end{table}


\paragraph{Gain on representative configurations}
The cache-bound case pairs a small testbed model whose fast decode
makes the workload prefill-dominated with a tight tool-call schedule;
\spear cuts mean TTFT by 35\% and tail end-to-end latency by 72\%. The
capacity-stressed case uses a loose schedule on a halved pool and
still delivers a 2.5\,pp hit-rate gain with a 6\% mean TTFT
reduction. The native-model case replaces the testbed with the
trajectory-native 30B under the tight schedule and confirms a positive
hit-rate gain of 1.6\,pp, ruling out a model-mismatch artifact.

\paragraph{Degradation under specific regimes}
Three traces not tabulated exhibit degraded or neutral outcomes.
On GAIA/Devstral under a halved pool and a tight schedule, tool gaps
fall below the decode latency, so the software idle timestamp inverts
the gap signal and hit rate drops by $6.6$\,pp. GAIA/Gemma4 has only
128 sessions, producing a noisy hazard LUT whose variance masks the
\spear advantage. SWE/Qwen3 is GPU-bound with negligible cache
pressure, matching the queue-bound SWE traces in
Tab.~\ref{tab:main}.

\paragraph{Efficacy boundary and hardware implication}
The degradation pattern reveals two efficacy boundaries of the
software patch. First, when the tool gap is shorter than the decode
window, a block-cache timestamp cannot resolve the idle onset and the
gap signal reverses, which is a fundamental limitation of any
software-level approximation that lacks a precise request-completion
event. Second, when prefill already saturates the GPU, a cache hit
saves tokens but does not shorten the queuing delay, so hit-rate gains
do not translate to latency gains. The first boundary is precisely what the event-driven hardware
interface of \unison eliminates, because a dedicated event FIFO stamps
gap onset and closure at cycle granularity and thereby removes the
software timing approximation that fails on short gaps. Note that the
\tide migration mechanism is not evaluated in this vLLM integration
because the production block cache exposes no tier-placement API, so
only the \spear ranking is exercised here.
